\Article{11}
Da die Bestimmung der plausibelsten Fehler von den Koeffizienten \(l,l',l''\) etc. unabhängig ist, stellt sie offenbar die bequemste Vorbereitung für jede beliebige Verwendung der Beobachtungen dar. Ferner ist ersichtlich, dass hierzu keine \textit{unbestimmte} Elimination oder Kenntnis der Koeffizienten \([\alpha\alpha],[\alpha\beta]\) etc. nötig ist; es ist lediglich erforderlich, die Hilfsgrößen \(A,B,C\) etc., die wir im Folgenden \textit{Korrelate} der Bedingungsgleichungen \(X=0,\ Y=0,\ Z=0\) etc. nennen werden, aus den Gleichungen (12) durch bestimmte Elimination zu ermitteln und in die Formeln (10) einzusetzen.

Obwohl diese Methode nichts zu wünschen übrig lässt, wenn nur die plausibelsten Werte der von den Beobachtungen abhängigen Größen gesucht sind, scheint es sich dennoch anders zu verhalten, wenn außerdem das Gewicht einer Bestimmung gewünscht wird. Denn hierzu scheint, je nachdem, welchen der vier oben angegebenen Ausdrücke man verwenden will, die Kenntnis der Größen \(L,L',L''\) etc. oder wenigstens der Größen \(x^0,y^0,z^0\) etc. erforderlich zu sein. Daher wird es nützlich sein, die Elimination genauer zu untersuchen; daraus wird sich uns auch ein leichterer Weg zur Ermittlung der Gewichte eröffnen.

\Article{12}
Der Zusammenhang der in dieser Untersuchung auftretenden Größen wird durch die Einführung der Funktion zweiten Grades in den Unbestimmten
\[
\begin{aligned}&[aa]xx+2[ab]xy+2[ac]xz+\text{etc.}\\&\qquad+[bb]yy+2[bc]yz+\text{etc.}+[cc]zz+\text{etc.}\end{aligned}
\]
die wir mit \(T\) bezeichnen wollen, beträchtlich verdeutlicht. Zunächst ist unmittelbar ersichtlich, dass diese Funktion gleich
\SourceLeaf{074}
\[
\begin{aligned}&\frac{(ax+by+cz+\text{etc.})^2}{p}+\frac{(a'x+b'y+c'z+\text{etc.})^2}{p'}\\&\qquad+\frac{(a''x+b''y+c''z+\text{etc.})^2}{p''}+\text{etc.}\quad(15)\end{aligned}
\]
wird.

Ferner ist klar, dass
\[
T=x\xi+y\eta+z\zeta+\text{etc.}\qquad(16)
\]
gilt; werden hierin wiederum \(x,y,z\) etc. mithilfe der Gleichungen (7) durch \(\xi,\eta,\zeta\) etc. ausgedrückt, so ist
\[
\begin{aligned}
T={}&[\alpha\alpha]\xi\xi+2[\alpha\beta]\xi\eta+2[\alpha\gamma]\xi\zeta+\text{etc.}\\&+[\beta\beta]\eta\eta+2[\beta\gamma]\eta\zeta+\text{etc.}\\
&+[\gamma\gamma]\zeta\zeta+\text{etc.}
\end{aligned}
\]

Die oben entwickelte Theorie enthält zwei Systeme bestimmter Werte der Größen \(x,y,z\) etc. sowie \(\xi,\eta,\zeta\) etc.; dem ersten, bei dem \(x=x^0,y=y^0,z=z^0\) etc. und \(\xi=-[al],\ \eta=-[bl],\ \zeta=-[cl]\) etc. gelten, entspricht der folgende Wert von \(T\):
\[
T=[ll]-\frac1P
\]
Dies wird entweder aus dem dritten Ausdruck für das Gewicht \(P\) im Vergleich mit der Gleichung (16) oder aus dem vierten unmittelbar ersichtlich; dem zweiten, bei dem \(x=A,y=B,z=C\) etc. sowie \(\xi=\mathfrak A,\ \eta=\mathfrak B,\ \zeta=\mathfrak C\) etc. gelten, entspricht der Wert \(T=S\), wie entweder aus den Formeln (10) und (15) oder aus (14) und (16) hervorgeht.

\clearpage
